\chapter*{Introduction}
\addcontentsline{toc}{chapter}{Introduction}

Why does physics repeatedly demand new mathematics?  A vibrating string first
asks only that possible motions can be added.  Quantum measurement then asks
for inner products and spectral projections.  Molecular spectra ask which
transformations preserve a system.  Electromagnetism asks how local data fit
together over spacetime.  Quantized Hall response asks which global features
survive deformation.  Open quantum systems ask that positivity remain valid
after arbitrary auxiliary systems are attached.  Anyons finally ask for a
language in which fusion and exchange are operations rather than afterthoughts.

This book follows those demands in the order
\[
\begin{gathered}
\text{superposition}\longrightarrow\text{symmetry}\longrightarrow
\text{fields}\\[-2pt]
\longrightarrow\text{topology}\longrightarrow\text{noncommutativity}
\longrightarrow\text{categorification}.
\end{gathered}
\]
The sequence is conceptual, not historical.  Physics decides why a structure
enters; mathematical prerequisites decide exactly where it can enter.

Each chapter begins with one physical question.  It then introduces the
smallest coordinate-free language that makes the question precise, proves one
central structural theorem, returns to the physical phenomenon, and ends with
exercises that connect the structure to models or experiments.  An
experimental observation is treated as data to be predicted or organized,
not as a proof of a mathematical formalism.  Every model states the
assumptions under which its conclusions are intended.

The companion volume \emph{Coordinate-free Mathematics} develops a broader
Algebra--Geometry--Analysis narrative.  The present book is not a rearranged
edition of it.  It has independent numbering, transitions, proofs, exercises,
and scope; the provenance ledger records the limited passages adapted from
the companion.
